% Источники: exam/materials/моделирование.pdf, раздел 24; https://github.com/HumsterProgrammer/iu9-notes/blob/main/8sem/Моделирование/Моделирование_лекция-05_26-02-26.md.
\section{Внешняя синхронизация имитационной модели.}

\textbf{Внешняя синхронизация} организует согласование компонент модели в системе дискретных событий. Она требуется, когда алгоритм функционирования компонент имитационной модели нужно делить на части операторами ожидания или остановки.

Используются операторы синхронизации:
\begin{itemize}
  \item немедленного перевода процесса в ждущее состояние;
  \item перевода процесса в ждущее состояние до выполнения некоторого условия.
\end{itemize}

Такими операторами алгоритм функционирования компонент имитационной модели делится на части.

Для транзактного и агрегатного способов внешняя синхронизация не требуется.

Важно, чтобы результаты функционирования некоторого кванта процесса появились в поле взаимодействия компонент сложной системы. Для этого используют рабочий массив, доступный только для данного процесса, и оператор условного ожидания некоторого момента времени. До наступления этого момента в рабочий массив переписываются показатели нового состояния всех компонент системы; при наступлении момента сформированный рабочий массив становится доступен другим компонентам для считывания обновленной информации.

\clearpage
